\section*{अध्याय \ref{ch:b2:retrosynthesis} --- पश्च-संश्लेषण और बहुचरणीय रणनीति}

\begin{solution}{exo:b2:retrosynthesis:1}\emergencystretch=3em
\begin{itemize}
  \item 2-फ़ेनिलप्रोपेन-2-ऑल: प्रोपेनोन + \ce{PhMgBr}, या ऐसीटोफ़ीनोन + \ce{CH3MgBr}।
  \item पेंटेन-3-ऑल: प्रोपेनल + \ce{CH3CH2MgBr}।
  \item 1-साइक्लोहेक्सिलएथेनॉल: एथेनल + साइक्लोहेक्सिलमैग्नीशियम ब्रोमाइड, या साइक्लोहेक्सेनकार्बैल्डिहाइड +
        \ce{CH3MgBr}।
  \item 2-मेथिलब्यूटेन-2-ऑल: प्रोपेनोन + \ce{CH3CH2MgBr}, या ब्यूटेनोन + \ce{CH3MgBr}।
  \item 1-फ़ेनिलएथेनॉल: बेन्ज़ैल्डिहाइड + \ce{CH3MgBr}, या एथेनल + \ce{PhMgBr}।
\end{itemize}
\end{solution}

\begin{solution}{exo:b2:retrosynthesis:2}
ग्राही \omterm{def:b2:retrosynthesis:disconnection}{सिंथॉन} $\mathrm{Ph\overset{+}{C}H{-}OH}$, तुल्य बेन्ज़ैल्डिहाइड; दाता \omterm{def:b2:retrosynthesis:disconnection}{सिंथॉन} \ce{CH3-},
तुल्य मेथिलमैग्नीशियम ब्रोमाइड (या मेथिललिथियम)। (दूसरा कटाव, \ce{Ph-} को दाता मानकर,
\ce{PhMgBr} और एथेनल प्रयुक्त करता है।)
\end{solution}

\begin{solution}{exo:b2:retrosynthesis:3}
$0.90 \times 0.85 \times 0.80 \times 0.95 \times 0.70 = 0.407$: लगभग 41~\%।
\end{solution}

\begin{solution}{exo:b2:retrosynthesis:4}
एस्टर को ऐल्कोहॉल तक अपचित करने के लिए आवश्यक \ce{LiAlH4} कीटोन का भी अपचयन कर देगा। कीटोन का
उसके चक्रीय ऐसीटैल के रूप में रक्षण कीजिए (एथेन-1,2-डाइऑल, अम्ल उत्प्रेरक, जल हटाते हुए); \ce{LiAlH4} से
एस्टर का अपचयन कीजिए; तनु जलीय अम्ल से ऐसीटैल हटाइए: 5-हाइड्रॉक्सीपेंटेन-2-ओन।
\end{solution}

\begin{solution}{exo:b2:retrosynthesis:5}
FGI: ब्यूटेन-1,3-डाइऑल $\Rightarrow$ 3-हाइड्रॉक्सीब्यूटेनल (ऐल्डिहाइड का अपचयन, \ce{NaBH4})। इस $\beta$-हाइड्रॉक्सी ऐल्डिहाइड का
$\alpha$--$\beta$ आबंध पर द्वि-समूह \omterm{def:b2:retrosynthesis:disconnection}{विच्छेदन}: एथेनल के दो अणु
(\omterm{def:b2:enolates-aldol:aldol}{ऐल्डोल})। आगे: एथेनल, तनु क्षारक, ठंडा; फिर \ce{NaBH4}।
\end{solution}

\begin{solution}{exo:b2:retrosynthesis:6}
साइक्लोहेक्सीन वलय को \omterm{def:b2:frontier-orbitals:diels-alder}{डील्स--ऐल्डर अभिक्रिया} में बने दो \ce{C-C} आबंधों पर काटा जाता है, जो
प्रतिस्थापी की ओर वलय द्वि-आबंध के सापेक्ष ऐलिलिक हैं: ब्यूटा-1,3-डाईन और ब्यूट-3-ईन-2-ओन,
ऐसीटिल समूह \omterm{def:b2:frontier-orbitals:diels-alder}{डाईनरागी} पर।
\end{solution}

\begin{solution}{exo:b2:retrosynthesis:7}
C1 (\ce{C=O}), C2=C3, फिर दोनों मेथिलों वाला C4 क्रमांकित करते हुए
(\cref{met:b2:conjugate-additions:robinson-disconnection}): \omterm{def:b2:enolates-aldol:aldol}{ऐल्डोल} कटाव C2--C3 C3 पर एक ऐल्डिहाइड
और एक मेथिल कीटोन छोड़ता है; माइकल कटाव C4--C5 2-मेथिलप्रोपेनल (C3 उसका कार्बोनिल, C4 उसका
$\alpha$ कार्बन) और ब्यूट-3-ईन-2-ओन छोड़ता है।
\end{solution}

\begin{solution}{exo:b2:retrosynthesis:8}
एक चरण: ब्यूटेनॉयल क्लोराइड और \ce{AlCl3} के साथ बेन्ज़ीन (\omterm{def:b2:aromatic-substitution:friedel-crafts}{फ़्रीडेल--क्राफ़्ट्स ऐसिलन}, कोई पुनर्विन्यास नहीं,
एकल ऐसिलन)। दो चरण: प्रोपिलमैग्नीशियम ब्रोमाइड के साथ बेन्ज़ैल्डिहाइड, फिर द्वितीयक
ऐल्कोहॉल का ऑक्सीकरण। फ़्रीडेल--क्राफ़्ट्स मार्ग छोटा है और उसके अभिकर्मक सस्ते, परंतु वह ऐलुमिनियम क्लोराइड का
पूरा तुल्यांक प्रयुक्त करता है, जो जलीय अपशिष्ट में पहुँचता है।
\end{solution}

\begin{solution}{exo:b2:retrosynthesis:9}
डाइक्लोरोमेथेन में पिरिडीनियम क्लोरोक्रोमेट, स्वर्न ऑक्सीकरण या डेस--मार्टिन पीरियोडिनेन
हेक्सेनल देते हैं; अम्लीकृत जलीय डाइक्रोमेट या परमैंगनेट हेक्सेनोइक अम्ल देते हैं, क्योंकि जल में
ऐल्डिहाइड अपना जलयोजित रूप बनाता है, जो फिर ऑक्सीकृत होता है (\cref{prop:b2:retrosynthesis:mild-oxidation})।
\end{solution}

\begin{solution}{exo:b2:retrosynthesis:10}
यदि ऐल्कोहॉल रक्षित हो तो ऐमीन बिना रक्षण के पहले अभिक्रिया करता है: ऐल्कोहॉल का सिलिल ईथर के रूप में रक्षण कीजिए
(यह ऐसिलन में बचा रहता है); ऐमीन का ऐसिलन कीजिए; फ़्लुओराइड से सिलिल समूह हटाइए; ऐल्कोहॉल का
ऐसिलन कीजिए। यदि ऐल्कोहॉल की अभिक्रिया के समय ऐमीन को मुक्त रखना होता, तो उसका
\textit{tert}-ब्यूटॉक्सीकार्बोनिल कार्बामेट (अम्ल से हटने वाला) के रूप में रक्षण किया जाता, जो सिलिल ईथर (फ़्लुओराइड से
हटने वाला) के लांबिक है: किसी को भी दूसरे को छुए बिना मुक्त किया जा सकता है।
\end{solution}

\begin{solution}{exo:b2:retrosynthesis:11}
\ce{CH3COCH2CH2COCH3} एक 1,4-डाइकार्बोनिल है: केंद्रीय \ce{C-C} आबंध काटिए। दाता
एथिल 3-ऑक्सोब्यूटेनोएट का \omterm{def:b2:enolates-aldol:enolate}{ईनोलेट} है (सोडियम एथॉक्साइड); ग्राही, किसी कीटोन का एक $\alpha$ कार्बन,
1-ब्रोमोप्रोपेन-2-ओन, \ce{BrCH2COCH3}, से मिलता है (ब्रोमाइड का द्वि-अणुक प्रतिस्थापन)। फिर एस्टर का जलीय जल-अपघटन
और तापन \ce{CO2} हटाते हैं: हेक्सेन-2,5-डाइओन।
\end{solution}

\begin{solution}{exo:b2:retrosynthesis:12}
पश्च-संश्लेषण: बगल में \ce{CH2} वाला मेथिल कीटोन ऐसीटोऐसीटिक एस्टर संश्लेषण का उत्पाद है;
C3 और C4 के बीच का आबंध काटिए: एथिल 3-ऑक्सोब्यूटेनोएट का \omterm{def:b2:enolates-aldol:enolate}{ईनोलेट} और ऐलिलिक हैलाइड
1-ब्रोमो-3-मेथिलब्यूट-2-ईन, \ce{(CH3)2C=CHCH2Br}। आगे: एथेनॉल में सोडियम एथॉक्साइड, फिर हैलाइड;
फिर जलीय क्षारक, अम्लीकरण और तापन, जो एस्टर का जल-अपघटन करते हैं और \ce{CO2} हटाते हैं
(\cref{prop:b2:enolates-aldol:decarboxylation}): 6-मेथिलहेप्ट-5-ईन-2-ओन।
\end{solution}

\begin{solution}{pb:b2:retrosynthesis:1}
\textbf{1.} \ce{HO-C6H4-CH2CH2COCH3}, \ce{OH} शृंखला के पैरा स्थान पर: एक फ़ीनॉल और एक कीटोन।
\textbf{2.} किसी \ce{C=C} का हाइड्रोजनन: \ce{ArCH2CH2CO} $\Rightarrow$ \ce{ArCH=CHCO}, एक
$\alpha,\beta$-असंतृप्त कीटोन।
\textbf{3.} (\textit{E})-4-(4-हाइड्रॉक्सीफ़ेनिल)ब्यूट-3-ईन-2-ओन, \ce{HOC6H4CH=CHCOCH3}।
\textbf{4.} एक $\alpha,\beta$-असंतृप्त कार्बोनिल: \omterm{def:b2:enolates-aldol:aldol}{ऐल्डोल संघनन}, \ce{C=C} पर कटाव।
\textbf{5.} 4-हाइड्रॉक्सीबेन्ज़ैल्डिहाइड और प्रोपेनोन।
\textbf{6.} ऐल्डिहाइड का कोई $\alpha$ हाइड्रोजन नहीं है और वह केवल इलेक्ट्रॉनरागी हो सकता है; आधिक्य में लिया गया प्रोपेनोन
ही अकेला \omterm{def:b2:enolates-aldol:enolate}{ईनोलेट} स्रोत है और स्वयं के बजाय अधिक अभिक्रियाशील ऐल्डिहाइड से अभिक्रिया करता है।
\textbf{7.} (1) 4-हाइड्रॉक्सीबेन्ज़ैल्डिहाइड, आधिक्य प्रोपेनोन, जलीय सोडियम हाइड्रॉक्साइड; अम्लीकरण। (2)
कमरे के ताप पर कार्बन पर पैलेडियम के ऊपर \ce{H2}, जब तक एक तुल्यांक ग्रहण न हो जाए।
\textbf{8.} \ce{CHO} समूह फ़ीनॉक्साइड को स्थायी करता है, अतः यह फ़ीनॉल कम-से-कम फ़ीनॉल जितना अम्लीय है
($\mathrm pK_a$ 9.99)। $\mathrm{pH} > 13$ पर $[\ce{ArO-}]/[\ce{ArOH}] > 10^{13-10} = 10^3$:
फ़ीनॉक्साइड।
\textbf{9.} \ce{O-} वलय से होकर कार्बोनिल कार्बन पर इलेक्ट्रॉन घनत्व धकेलता है (उसके पैरा स्थित एक
अनुनादी दाता): ऐल्डिहाइड कम इलेक्ट्रॉनरागी है, और संघनन धीमा।
\textbf{10.} बेन्ज़िल ईथर: \ce{C=C} का अपचयन करने वाला पैलेडियम पर हाइड्रोजनन
बेन्ज़िल ईथर का भी विदलन कर देता है (हाइड्रोजन-अपघटन), और उसी चरण में फ़ीनॉल मुक्त हो जाता है।
\textbf{11.} तीन: बेन्ज़िलन, \omterm{def:b2:enolates-aldol:aldol}{ऐल्डोल संघनन}, विरक्षण के साथ हाइड्रोजनन।
\textbf{12.} नहीं: मृदु परिस्थितियों में (कमरे का ताप, पैलेडियम) \omterm{def:b2:huckel:aromatic}{ऐरोमैटिक} वलय का अपचयन नहीं होता।
\textbf{13.} रक्षण नहीं: फ़ीनॉक्साइड केवल संघनन को धीमा करता है, जिसकी क्षतिपूर्ति अधिक समय या तापन से
हो जाती है; एक चरण की बचत दर के लाभ से अधिक मूल्यवान है।
\textbf{14.} 4-आयोडोफ़ीनॉल और ब्यूट-3-ईन-2-ओन।
\textbf{15.} \ce{HOC6H4I + CH2=CHCOCH3 + Et3N -> HOC6H4CH=CHCOCH3 + Et3NH+ + I-}, पैलेडियम उत्प्रेरक।
\textbf{16.} अंतस्थ \ce{CH2} पर, जो $\beta$ कार्बन है: ऐरिल समूह कम बाधित सिरे पर निविष्ट होता है,
और फिर $\beta$-हाइड्राइड विलोपन संयुग्मित (\textit{E})-ईनोन देता है।
\textbf{17.} \ce{Pd(0)} पर ऐरिल आयोडाइड का \omterm{def:b2:catalytic-cycles:oxidative-addition}{ऑक्सीकारी योग}, ऐल्कीन का उपसहसंयोजन और निवेशन,
$\beta$-हाइड्राइड विलोपन, और क्षारक द्वारा \ce{Pd(0)} का पुनर्जनन, जो \ce{HI} ले लेता है
(\cref{ch:b2:catalytic-cycles})।
\textbf{18.} $162.188/(220.005 + 70.091 + 101.193) = 162.188/391.289 \approx 41.4~\%$।
\textbf{19.} \ce{C7H6O2 + C3H6O -> C10H10O2 + H2O}: $162.188/(122.123 + 58.080) = 162.188/180.203
\approx 90.0~\%$।
\textbf{20.} 100~\%: \ce{H2} और ईनोन का हर परमाणु उत्पाद में पहुँचता है।
\textbf{21.} $164.204/(122.123 + 58.080 + 2.016) = 164.204/182.219 \approx 90.1~\%$।
\textbf{22.} \ce{C=C}: संयुग्मित और अबाधित, यह पैलेडियम पर बँधता है और कीटोन \ce{C=O} से कहीं
तेज़ी से हाइड्रोजनित होता है; हाइड्रोजन के एक तुल्यांक के बाद रोक देने से कीटोन बचा रहता है।
\textbf{23.} \omterm{def:b2:enolates-aldol:aldol}{ऐल्डोल} मार्ग: सस्ता ऐल्डिहाइड, प्रोपेनोन और सोडियम हाइड्रॉक्साइड, उपोत्पाद जल। हेक
मार्ग: एक ऐरिल आयोडाइड, एक पैलेडियम उत्प्रेरक, और ब्यूट-3-ईन-2-ओन, जो अत्यंत विषैला है।
\textbf{24.} \omterm{def:b2:enolates-aldol:aldol}{ऐल्डोल} मार्ग: दो चरण, सस्ते और सुरक्षित अभिकर्मक, 90~\% से अधिक परमाणु मितव्ययिता।
\textbf{25.} यह उसे 0.90 से गुणा कर देता: $0.782 \times 0.90 \approx 70~\%$।
\textbf{26.} समग्र $0.85 \times 0.92 = \boldsymbol{\approx 78~\%}$।
\end{solution}
